\chapter*{Sources and further reading}
\addcontentsline{toc}{chapter}{Sources and further reading}
Access checked 17 October 2026. The reference builds on the actual Chapter 31--33 native, pool and checkpoint contracts rather than on a detached recurrence or slot dictionary. Its tests, tables, code listings and illustrations are bound to those sources.

The USENIX Orca presentation-page abstract was inspected for iteration-level
scheduling, not for a reproduced serving implementation. NVIDIA's live
CUDA Best Practices Guide 13.4 was inspected on profiling, verification and
coalesced accesses. No external benchmark is reproduced or transferred to
DOGMA. Mathematical derivations and arithmetic models in this chapter are
original and explicitly restricted to the declared recurrence.

The CPU reference is not a fused device kernel, concurrent service, durable
work journal, authenticated ownership system or measured genomic model.
Local technical and visual review does not replace independent review.
\begin{thebibliography}{9}
\bibitem{orca} Gyeong-In Yu, Joo Seong Jeong, Geon-Woo Kim, Soojeong Kim
and Byung-Gon Chun. \emph{Orca: A Distributed Serving System for
Transformer-Based Generative Models}. OSDI 2022, pp.\ 521--538.
Presentation-page abstract inspected.
\url{https://www.usenix.org/conference/osdi22/presentation/yu}.
\bibitem{cuda} NVIDIA. \emph{CUDA C++ Best Practices Guide}, version 13.4,
sections on profiling, reference comparison, numerical accuracy and
coalesced access to global memory. Live documentation inspected; no
hardware performance validation performed.
\url{https://docs.nvidia.com/cuda/cuda-c-best-practices-guide/index.html}.
\end{thebibliography}
